\section*{AI-use statement}

This paper was drafted by a crew of AI research-assistant agents
(Free-Research-Crew, protocol v0.1) under the direction of the human author, Bryce Wang.
He approved the research question (2026-11-19, 01:17 UTC), the design (2026-11-19, 07:34 UTC)
and the results (2026-11-19, 13:05 UTC) at three human gates, each in chat after reading the
Reviewer's verdict and kill argument. The manuscript itself is subject to a fourth, the
submission gate. He takes responsibility for the paper's content.

\begin{itemize}
  \item \textbf{Roles and models.} Every crew role ran on Anthropic Claude through Claude
        Code. The PI role was the main Claude Code session. The Scout, Designer, Data RA,
        Econometrician, Figures RA, Reviewer, Writer and Replicator ran as agents launched
        from that session; the review records name the model as \texttt{claude-opus-5-5}.
        The PI's brief to each agent is archived in \texttt{briefs/}, from the Scout's first
        brief (2026-11-18, 16:28 UTC) to this Writer's brief (2026-11-19, 13:44 UTC).
  \item \textbf{Review.} The Reviewer's independence level was \emph{fresh-context}: a new
        context in every round, with the same model family as the other roles. The human
        author chose this over a cross-family review. The Reviewer ran 4 rounds on the
        research question (\emph{block}, \emph{revise}, \emph{revise}, \emph{pass}), 2 on the
        design (\emph{block}, with one critical finding, then \emph{pass}), 1 on the code
        (\emph{pass} with four minor findings) and 2 on the results (\emph{revise}, then
        \emph{pass}). No separate literature round was held. The review of this paper has not
        yet taken place. The question (round~4), design (round~2) and results (round~2)
        stages each ended with a \emph{pass} and minor findings still open, which were
        answered afterwards: the question wording by the PI, checked by the Reviewer at the
        first design round; the design wording before the design gate; and the claim texts
        and run notes after the results gate. The code review's four minor findings were fixed
        and the analysis then ran without a second code-review round, at the PI's instruction;
        the Reviewer checked the fixes at the first results round
        (Section~\ref{sec:deviations}, disclosure~1). A post-results arithmetic run was not
        code-reviewed (disclosure~2).
  \item \textbf{Findings accepted over the Reviewer's objection.} None. No gate was approved
        with an open critical finding, and no gate note records an override. One question-stage
        finding (that the original H2 was close to unfalsifiable) was conceded by narrowing H2
        to a descriptive comparison.
  \item \textbf{Evidence.} Every result number in the text and tables is a macro generated
        from the metrics of a logged run, or is printed by a logged table script from run
        metrics. Numbers typed by hand are design facts fixed in the plan before any analysis
        run (sample counts, folds, seeds, the propensity clip, the placebo bound, basis sizes,
        the outcome-free overlap diagnostics), descriptive covariate statistics from the
        outcome-free data-preparation and diagnostic runs (Table~\ref{tab:summary}), facts from
        the cited sources as the Scout's notes record them, and process facts such as dates and
        round counts. Every citation was checked against its source by the Scout: of the 17
        sources, 14 were read in full text and 3 from the abstract only, and each is cited for
        no more than was read. The full workspace is public at
        \url{https://github.com/brycewang-stanford/Free-Research-Crew/tree/main/gallery/ss-lalonde-observational}.
        It contains the plan with its deviations and disclosures, the run ledger and notes, the
        code, the data manifest (the data files themselves are fetched, not committed), every
        review, the archived briefs and the source notes.
  \item \textbf{Replication.} Passed. The Replicator's clean rerun, recorded in
        \texttt{replication.json} (2026-11-19, 13:59 UTC), copied the workspace, deleted every
        recorded output and re-executed all 19 logged runs, from the data download to the
        figures and the main table, in a fresh environment built from
        \texttt{code/requirements.txt} alone. The data were fetched again rather than reused.
        All 19 runs reproduced every recorded metric within the protocol tolerance (relative
        $10^{-6}$, absolute $10^{-9}$). No code, data or tolerance was changed for it. The
        Replicator's account is in \texttt{REPLICATION.md}.
\end{itemize}
